\documentclass[11pt]{article}

\usepackage[margin=1in]{geometry}
\usepackage{amsmath,amssymb,amsthm}
\usepackage{enumitem}
\usepackage{hyperref}

\newtheorem{proposition}{Proposition}
\newtheorem{lemma}{Lemma}

\newcommand{\calR}{\mathcal{R}}
\newcommand{\calP}{\mathcal{P}}
\newcommand{\calQ}{\mathcal{Q}}
\newcommand{\calW}{\mathcal{W}}
\newcommand{\calX}{\mathcal{X}}



\begin{document}

\begin{itemize}
\item Define
\[
  \sigma(n):=\sum_{d\mid n}d
\]
be the divisor-sum function.
\item Define
\[
  S(t):=\sum_{m\geq0}\sigma(m+1)(-t)^m
       =1-3t+4t^2-7t^3+6t^4-12t^5+\cdots.
\]
\item  Define coefficients $c_m$ by
\[
  C(t):=\sum_{m\geq0}c_m t^m:=\frac1{S(t)}.
\]
\item Define
\[
  P(t):=\frac{1+t}{1-t}=1+2t+2t^2+2t^3+\cdots
\]
\item Define
\[
  Q(t):=1-P(t)S(t)=\sum_{n\geq1}q_n t^n.
\]

\item For $m\geq1$, define $\calR_m$ be the set of triples $(a,b,i)$ such that
\[
  a,b\geq1,
  \qquad ab=m,
  \qquad 1\leq i\leq a.
\]

\item For $N\geq2$, define a finite colored set
\[
  \calX_N:=
  \{(a,b,i,\varepsilon):ab\leq N,
       \ 1\leq i\leq a,
       \ \varepsilon\in E_N(ab)\},
\]
where
\[
  E_N(m):=
  \begin{cases}
    \{0,1\},&m<N,\\
    \{0\},&m=N.
  \end{cases}
\]

\end{itemize}


Thus, rectangles of area $<N$ occur in two colors, and rectangles of area $N$ occur only in color $0$.  Call an object {\it positive} if its area $ab$ is even, and {\it negative} if its area is odd.  Let $\calX_N^+$ and $\calX_N^-$ denote the two parts.  The unusual coloring rule at the boundary is exactly what matches the final, un-doubled term in the coefficient computation below.

\begin{lemma}\label{lem:q-signed-count}
For $n\geq1$, with $N=n+1$, we have
\[
  q_n=|\calX_N^+|-|\calX_N^-|.
\]
\end{lemma}

\begin{proof}
Since $P(t)=1+2t+2t^2+\cdots$, we have
\[
  [t^n]P(t)S(t)
  =(-1)^n\sigma(n+1)+2\sum_{m=0}^{n-1}(-1)^m\sigma(m+1).
\]
Putting $N=n+1$ and replacing $m+1$ by $r$ gives
\[
  q_n=-[t^n]P(t)S(t)
      =(-1)^N\sigma(N)+2\sum_{r=1}^{N-1}(-1)^r\sigma(r).
\]
This is exactly the signed count of $\calX_N$, because each area $r<N$ has two colored copies of $\calR_r$, area $N$ has one copy, and the sign is $+1$ for even area and $-1$ for odd area.
\end{proof}

\begin{lemma}\label{lem:injection}
For every $N\geq2$, there is an explicit injection
\[
  \Phi_N:\calX_N^-\hookrightarrow\calX_N^+.
\]
Consequently $q_n\geq0$ for every $n\geq1$.
\end{lemma}

\begin{proof}
Let $(a,b,i,\varepsilon)\in\calX_N^-$.  Then $ab$ is odd, so $a$ and $b$ are odd.

\medskip
\noindent\textbf{Case 1: $b>1$.}
Set
\[
  \Phi_N(a,b,i,\varepsilon):=(a,b-1,i,\varepsilon).
\]
The new width is positive and even, the new area is $a(b-1)<ab\leq N$, and both colors are allowed.

\medskip
\noindent\textbf{Case 2: $b=1$ and $a<N$.}
Here both colors are available.  If $\varepsilon=0$, set
\[
  \Phi_N(a,1,i,0):=(a+1,1,i,0).
\]
If $\varepsilon=1$ and $i=1$, set
\[
  \Phi_N(a,1,1,1):=(a+1,1,a+1,0).
\]
If $\varepsilon=1$ and $i>1$, set
\[
  \Phi_N(a,1,i,1):=(a-1,1,i-1,1).
\]
These are positive columns with valid marked rows.  The first two have height at most $N$, and use color $0$ if the boundary is reached; the third has area $<N$, so color $1$ is allowed.

\medskip
\noindent\textbf{Case 3: $b=1$ and $a=N$.}
This only occurs for odd $N$, and only color $0$ is available.  If $i>1$, set
\[
  \Phi_N(N,1,i,0):=(N-1,1,i-1,1),
\]
and if $i=1$, set
\[
  \Phi_N(N,1,1,0):=(1,N-1,1,1).
\]
Both images have even area $N-1<N$, so color $1$ is allowed.

We now check injectivity.  Images from Case 1 are reversible by increasing the width by $1$; this recovers the original odd width and preserves the marked row and color.  The exceptional image $(1,N-1,1,1)$ from Case 3 cannot be a Case 1 image, because its only possible Case 1 preimage would be $(1,N,1,1)$, and color $1$ is forbidden at area $N$.

It remains to separate the column images.  A color $0$ column of even height $e$ has a unique source: rows $1,\ldots,e-1$ come from $(e-1,1,i,0)$, while row $e$ comes from $(e-1,1,1,1)$.  A color $1$ column of even height $e<N-1$ must come from $(e+1,1,i+1,1)$.  A color $1$ column of height $N-1$ must instead come from the boundary rule $(N,1,i+1,0)$.  These alternatives are disjoint and each gives at most one preimage.  Hence no two negative objects have the same image.  Therefore $\Phi_N$ is injective, and Lemma~\ref{lem:q-signed-count} gives $q_n\geq0$.
\end{proof}


For $r\geq0$, let $\calP_r$ be any set with
\[
  |\calP_r|=
  \begin{cases}
    1,&r=0,\\
    2,&r\geq1.
  \end{cases}
\]
These are the $P$-objects.

For $n\geq1$, put $N=n+1$ and define
\[
  \calQ_n:=\calX_N^+\setminus\Phi_N(\calX_N^-).
\]
Thus $\calQ_n$ is the set of unmatched positive row-marked rectangles after the injection.  By Lemmas~\ref{lem:q-signed-count} and~\ref{lem:injection},
\[
  |\calQ_n|=q_n.
\]

For $m\geq0$, let $\calW_m$ be the set of all finite assemblies
\[
  (U;A_1,A_2,\ldots,A_\ell)
\]
such that $\ell\geq0$, $U\in\calP_r$ for some $r\geq0$, each $A_j\in\calQ_{n_j}$ for some $n_j\geq1$, and
\[
  r+n_1+\cdots+n_\ell=m.
\]
The order of the $A_j$ matters.  Thus the $A_j$ form a word, not a set or multiset.

\begin{proposition}\label{prop:main}
For every $m\geq0$, we have
\[
  c_m=|\calW_m|.
\]
In particular, $c_m>0$ for all $m\geq0$.
\end{proposition}

\begin{proof}
The $P$-objects have generating function $P(t)$, while the $Q$-objects have generating function $Q(t)$.  A finite ordered word of $Q$-objects has generating function
\[
  1+Q(t)+Q(t)^2+\cdots=\frac1{1-Q(t)},
\]
since $Q(0)=0$.  Therefore the assemblies $\calW_m$ are generated by
\[
  P(t)\frac1{1-Q(t)}
  =P(t)\frac1{P(t)S(t)}
  =\frac1{S(t)}
  =C(t).
\]
Comparing coefficients gives $c_m=|\calW_m|$.

The set $\calW_m$ is nonempty for every $m$: for $m=0$, take the unique object of $\calP_0$ and the empty word; for $m\geq1$, take an object of $\calP_m$ and the empty word.  Hence $c_m>0$.
\end{proof}

\end{document}
